\documentclass[10pt,a4paper]{amsart}
\usepackage[margin=19mm]{geometry}
\usepackage{amsmath,amssymb,mathtools}
\usepackage[scr=rsfs,cal=euler]{mathalfa}
\usepackage{xfrac}
\usepackage[unicode,hidelinks,bookmarks=false]{hyperref}
\newcommand{\ie}{i.e.\ }
\newcommand{\cf}{cf.\ }
\newcommand{\resp}{respectively\ }
\newcommand{\Subsection}{\subsection}
\providecommand{\nobreakdash}{\nobreak\hbox{-}}
\setlength{\emergencystretch}{2em}
\multlinegap0pt
\usepackage{bm}
\usepackage{bbm}
\usepackage{dsfont}
\usepackage{xspace}
\usepackage{cancel}
\usepackage{mathtools}
\usepackage{amsthm}
\makeatletter
\@ifundefined{theorem}{\newtheorem{theorem}{Theorem}}{}
\makeatother
\newcommand{\CC}{\mathbb{C}}
\title[Grassmann and Flag Varieties in Linear Algebra, Optimization]{Grassmann and Flag Varieties in Linear Algebra, Optimization, and Statistics: An Algebraic Perspective}
\author{Hannah Friedman, Serkan Ho\c{s}ten}
\date{Source: arXiv:2505.15969v2 (CC BY 4.0)}
\begin{document}
\maketitle

\noindent\emph{Source and licence.} Grassmann and Flag Varieties in Linear Algebra, Optimization, and Statistics: An Algebraic Perspective. Version arXiv:2505.15969v2 (CC BY 4.0), distributed under the Creative Commons Attribution 4.0 International licence, as recorded in the arXiv licence metadata for this exact version and shown on the abstract page https://arxiv.org/abs/2505.15969v2. Full rights notice: licenses/arxiv-2505.15969-CC-BY-4.0.txt.\par\smallskip

\noindent\textit{Editorial scope.} Counted unit: the Theorem whose source label is \texttt{thm:canonical} in the article (source0.tex lines 810--814, source label \texttt{thm:canonical}), delivered together with the article's own complete original proof of it (source0.tex lines 815--865, one \texttt{proof} environment of 51 lines). No mathematics is written, repaired or re-derived: statement and proof are the article's own text line for line. Required context retained uncounted: eqn:ccabrokendown (lines 803--808). Every definition, display and label of those lines is retained unchanged. Omitted from this package and not redistributed: 1--802, 809--809, 866--1043. Nothing inside a retained excerpt is omitted or altered: every retained body is delivered exactly as the article's lines are written, so no omission receipt of any kind accompanies this package. Citations: the retained excerpts carry no citation command at all, so no bibliography entry of the article is transcribed and no reference is redistributed. Reference adaptation: none was needed and none was made; every reference inside the delivered unit resolves inside it, so no unresolved reference or citation is produced. Printed slips kept literal: no printed word, symbol, hypothesis or number of the counted statement or of its proof was altered. Layout and class: the article's own class is \texttt{extarticle} with packages including amsmath, amsthm, amssymb, color, hyperref, graphicx, caption, mathtools, enumerate, verbatim, tikz, tikz-cd, tikz-3dplot, algorithm; the wrapper is the lane's pinned \texttt{amsart} editorial wrapper at 10pt on A4, which restates the theorem-like environments the retained text needs and adds the article's own macro definitions that the retained text needs (\texttt{\textbackslash CC}), with extra wrapper packages (bm, bbm, dsfont, xspace, cancel, mathtools). The delivered numbering is the unit's own. Assets: no retained span contains a figure environment, a graphics-inclusion command, a PDF-page inclusion, a table or a TikZ picture; no image file is copied into this package, the package declares no asset dependency and no image file is redistributed with it. Licence: version arXiv:2505.15969v2 of this article is distributed under the Creative Commons Attribution 4.0 International licence, as recorded in the arXiv licence metadata for this exact version and shown on the abstract page https://arxiv.org/abs/2505.15969v2. A full rights notice is delivered at licenses/arxiv-2505.15969-CC-BY-4.0.txt. Characters: every retained excerpt is pure ASCII, so the delivered engine, XeLaTeX with its default Latin Modern text font, reports no missing character.
\begin{align}\label{eqn:ccabrokendown}
    \textrm{maximize } &u^T A v\\\nonumber
    \textrm{subject to } &u^Tu = v^Tv = 1\\\nonumber
    &u^Tu_j = u^TAv_j = 0, \,\, j=1, \ldots, k-1\\\nonumber
    &v^TA^Tu_j = v^Tv_j = 0,  \,\, j =1, \ldots, k-1.
\end{align}

\begin{theorem} \label{thm:canonical}
The critical points of the canonical correlation analysis problem are the pairs $(U, V) \in \CC^{p \times k} \times \CC^{q \times k}$
where the columns of $U$ are the left singular unit vectors of $A$, the columns of $V$ are the right singular unit vectors of $A$, and  corresponding columns have the same singular value.
There are $\binom{\min(p, q)}{k} k! 2^k $ critical points.
\end{theorem}

\begin{proof}
    The count follows from the description of the critical points: choose $k$ unit left and right singular vectors of $A$, permute the corresponding vectors simultaneously, and flip  signs of them as desired. 
    We will prove by induction that for each $i$, there exist $\lambda_i, \eta_i$ such that $Av_i = \lambda_i u_i$ and $A^T u_i = \eta_i v_i$ if and only if the pair $(u_i, v_i)$ is a critical point of \eqref{eqn:ccabrokendown}. 
    
    We proceed by induction on $i$.
    When $i = 1$, the optimization problem simplifies to 
    $\max_{u^Tu = v^Tv = 1} \,\, v^TAu.$
    By computing the transpose of the Jacobian, we find that the critical points of this problem are characterized by the leftmost column of the matrix
    \begin{align*}
        \begin{pmatrix}
            Av & u & 0\\
            A^T u & 0  & v
        \end{pmatrix}
    \end{align*}
    being in the span of the rest of the columns. 
    We therefore have that $(u,v)$ is a critical point if and only if there exist $\lambda_1, \eta_1$ such that $Av = \lambda_1 u$ and $A^Tu = \eta_1 v$. 

    Suppose now that for $j = 1, \ldots, i - 1$, we have that $Av_j = \lambda_j u_j$ and $A^T u_j = \eta_j v_j$. 
    The transpose of the augmented Jacobian matrix of \eqref{eqn:ccabrokendown} is 
    \begin{align*}
    &\left (
     \begin{array}{ccccccccccccccc}
            Av_i & u_1 & \cdots & u_i & 0 & \cdots & 0 & Av_1 & \cdots & Av_{i-1}& 0 & \cdots & 0\\
            A^Tu_i & 0 & \cdots & 0 & v_1 & \cdots & v_i & 0 & \cdots & 0 & A^Tu_1 & \cdots & A^Tu_{i-1}
        \end{array}
        \right)\\
        = &   \left (
     \begin{array}{ccccccccccccccc}
            Av_i & u_1 & \cdots & u_i & 0 & \cdots & 0 & \lambda_1 u_1 & \cdots & \lambda_{i-1}u_{i-1}& 0 & \cdots & 0\\
            A^Tu_i & 0 & \cdots & 0 & v_1 & \cdots & v_i & 0 & \cdots & 0 & \eta_1v_1 & \cdots & \eta_{i-1}v_{i-1}
        \end{array}
        \right).
     \end{align*}
    Proving that this matrix drops rank is equivalent to proving that the matrix
    \begin{align*}
         \begin{pmatrix}
            Av_i  & u_1 & \cdots & u_i & 0 & \cdots & 0\\
            A^T u_i & 0 & \cdots & 0 & v_1 & \cdots & v_i
            \end{pmatrix}
    \end{align*}
    drops rank. 
    It is clear that the matrix drops rank if $Av_i \in {\rm span}(u_i)$ and $A^Tu_i \in {\rm span}(v_i)$.
    Conversely, suppose that
    \begin{align*}
        Av_i = \alpha_1 u_1 + \cdots + \alpha_i u_i,\\
        A^T u_i = \beta_1 v_1 + \cdots + \beta_i v_i.
    \end{align*}
    Multiplying the first equation by $u_j^T$ for $j = 1, \ldots, i-1$ gives $0 = u_j^T A v_i = \alpha_j$. 
    By symmetry, $\alpha_j = \beta_j = 0$ for $j =1, \ldots, i-1$. 
    Thus $Av_i = \alpha_i u_i$ and $A^T u_i = \beta_i v_i$. 
\end{proof}

\end{document}
